% ============ Solutions and hints for exercises (Part 1: §1-§3) ============
% 原书 pp.287-288; 原书对 §1 的 Ex.1.4 与 §3 的 Ex.3.3 未给解答, 编号按原书跳号.

\noindent\textit{(Please be sure to try each exercise on your own before looking at
the solution)}

\section*{\S\ref{sec:1}}

\noindent\textbf{1.1.}\quad If $ab = 1 - x$ with $x^n = 0$ then $ab(1 + x + \cdots + x^{n-1}) = 1$.

\noindent\textbf{1.2.}\quad Set $e_i = (0, \ldots, 1, \ldots, 0) \in A_1 \times \cdots \times A_n$ (with $1$ in the $i$th place); then since
$e_ie_j = 0$ for $i \neq j$, any prime ideal $p$ of $A_1 \times \cdots \times A_n$ must contain all but one
of the $e_i$.

\noindent\textbf{1.3.}\quad (a) Use the fact that $\rad(A) = \{x \in A \mid 1 + ax \text{ is a unit of } A, \forall a \in A\}$.
Counter-example: $A = \Z$, $B = \Z/(4)$; then $\rad(A) = (0)$, $\rad(B) = 2B$.

\smallskip
\noindent\hspace*{1.5em}(b) Let $m_1, \ldots, m_r$ be the maximal ideals of $A$ and $I = \Ker f$. Suppose
$I \subset m_i$ for $1 \le i \le s$ and $I \nsubseteq m_i$ for $s < i \le r$; then the maximal ideals of $B$
are $f(m_i)$ for $1 \le i \le s$, and $f(m_i) = B$ for $i > s$. Now $\rad(A) = m_1 \cdots m_r$,
hence $f(\rad(A)) = f(m_1) \cdots f(m_r) = f(m_1) \cdots f(m_s) = \rad(B)$.

\noindent\textbf{1.5.}\quad The first half is easy; for the second, use Zorn's lemma.

\noindent\textbf{1.6.}\quad We can assume that there are no inclusions among $P_1, \ldots, P_r$. When $r = 2$,
take $x \in I - P_1$, $y \in I - P_2$; then one of $x$, $y$, $x + y$ is not in $P_1$ or $P_2$. When
$r > 2$, we can take $x \in I - (P_1 \cup \cdots \cup P_{r-1})$ by induction. Also, since $P_r$ is
prime, $P_r \not\supset IP_1 \cdots P_{r-1}$, so take $y \in IP_1 \cdots P_{r-1} - P_r$; then either $x$ or $x + y$
satisfies the condition.

\section*{\S\ref{sec:2}}

\noindent\textbf{2.1.}\quad By NAK there is an $e \in I$ such that $(1 - e)I = 0$. One sees easily that then
$I = Ie = Ae$ and $e^2 = e$.

\noindent\textbf{2.2.}\quad If $x \in \ann(M/IM)$ then $xM \subset IM$, so that by Theorem~\hyperref[thm:2.1]{1} there exists $y \in I$
such that $(x^n + y)M = 0$.

\noindent\textbf{2.3.}\quad $(M + N)/N \cong M/(M \cap N)$ shows that $M$ is finite, and similarly for $N$.

\noindent\textbf{2.4.}\quad (a) If $M \cong A^n$ and $P$ is a maximal ideal of $A$ with $k = A/P$ then $M/PM = k^n$;
for a field the result is well-known.

\smallskip
\noindent\hspace*{1.5em}(b) The first part is easy by the theory of determinants; the second half
comes from the fact that $A^n$ has $n$ linearly independent elements, but any
$n + 1$ elements are linearly dependent.

\smallskip
\noindent\hspace*{1.5em}(c) Use Theorem~\hyperref[thm:2.3]{3}, (iii).

\noindent\textbf{2.5.}\quad (a) If $F$ and $F'$ are free modules and $\alpha\colon F \to L$, $\beta\colon F' \to N$ are surjections
then there is a map $\mu$ making
\[
  \begin{tikzcd}[column sep=large, row sep=large]
    0 \arrow[r] & F \arrow[r] \arrow[d, "\alpha"] & F \oplus F' \arrow[r] \arrow[d, "\mu"] & F' \arrow[r] \arrow[d, "\beta"] & 0 \\
    0 \arrow[r] & L \arrow[r] & M \arrow[r] & N \arrow[r] & 0
  \end{tikzcd}
\]
commute. The assertion follows from this and the snake lemma
(Appendix B).

\smallskip
\noindent\hspace*{1.5em}(b) can be proved similarly.

\section*{\S\ref{sec:3}}

\noindent\textbf{3.1.}\quad Use the fact that $A$ is isomorphic to a submodule of $(A/I_1) \oplus \cdots \oplus (A/I_n)$.

\noindent\textbf{3.2.}\quad Use the previous question.

\noindent\textbf{3.4.}\quad If $II^{-1} = A$ then there exist $x_i \in I$ and $y_i \in I^{-1}$ such that $\sum x_iy_i = 1$; then it
follows easily that $I = \sum x_iA$.

\noindent\textbf{3.5.}\quad If $J$ is a fractional ideal generated by $a_1/b_1, \ldots, a_n/b_n$ with $a_i$ and
$b_i$ coprime, then $J^{-1}$ is the principal fractional ideal generated by $u/v$, where
$v = \lcm(b_1, \ldots, b_n)$ and $u = \hcf(a_1, \ldots, a_n)$.

\noindent\textbf{3.6.}\quad $\Ker(\varphi^n)$, $n = 1, 2, \ldots$, is an ascending chain of ideals of $A$.

\noindent\textbf{3.7.}\quad Choose an ideal $I$ of $A$ which is not finitely generated, and set $M = A/I$;
then by Theorem~\ref{thm:2.6}, $M$ cannot be of finite presentation.
